% !PRINTED 112
\weberpage{112}

\begin{center}
\textbf{\S\ 34. \gesperrt{The functions $\eta(\omega)$, $f(\omega)$, $f_1(\omega)$, $f_2(\omega)$.}}
\end{center}

The functions $\eta(\omega)$, $f(\omega)$, $f_1(\omega)$, $f_2(\omega)$ have already
been mentioned in § 24, where, in the representation of the $\vartheta$-functions by
infinite products, they presented themselves almost of their own accord as
single-valued functions of $\omega$, whereas they did not arise in the second
representation by infinite series. Connected with this is a remarkable circumstance:
many of our formulae, for example those belonging to the transformation of the second
order, (17), § 32, or formula (23), § 32, can easily be verified by means of the
infinite products, but with difficulty or not at all by means of the infinite series.
For that reason it was of interest to us to derive these results from the theory of
transformation without making use of either of these expressions, and in the same way
the functions $\eta(\omega)$, $f(\omega)$, $f_1(\omega)$, $f_2(\omega)$ and their
fundamental properties shall now likewise be obtained from this source.

Formula (6) of the preceding section gives for $n=3$:
\[
\sqrt{3}\,\vartheta_{11}(3u,\,3\omega)\,\vartheta_{00}\,\vartheta_{10}\,\vartheta_{01}
= 2\,\vartheta_{11}\!\left(\tfrac{1}{3}\right)\vartheta_{11}(u)\,
\vartheta_{11}\!\left(\tfrac{1}{3}+u\right)\vartheta_{11}\!\left(\tfrac{1}{3}-u\right),
\]
and if one differentiates and then sets $u=0$, one obtains by (4):
\[
3\sqrt{3}\,\vartheta'_{11}(0,\,3\omega) = 2\pi\,\vartheta_{11}\!\left(\tfrac{1}{3}\right)^3,
\]
or, replacing $\omega$ by $\dfrac{\omega}{3}$,
\[
3\sqrt{3}\,\vartheta'_{11} = 2\pi\left[\vartheta_{11}\!\left(\tfrac{1}{3},\,\tfrac{\omega}{3}\right)\right]^3.
\]
If one therefore puts
\weq{1}{\eta(\omega) = \frac{1}{\sqrt{3}}\,\vartheta_{11}\!\left(\frac{1}{3},\,\frac{\omega}{3}\right),}
then it follows, in the notation agreeing with § 24 (9):
\weq{2}{\vartheta'_{11} = \pi\,\vartheta_{00}\,\vartheta_{01}\,\vartheta_{10} = 2\pi\,\eta(\omega)^3}
and from § 25, (4) one obtains for $\eta(\omega)$ the series development
\weq{3}{\eta(\omega) = q^{\frac{1}{12}} \sum_{-\infty,\,\infty}^{\nu} (-1)^{\nu}\, q^{3\nu^2+\nu}.}

The function $\eta(\omega)$ is \gesperrt{real} for a purely imaginary $\omega$ (real $q$), and for
an infinitely large $\omega$ (i.\,e. vanishing $q$) we have
\[
q^{-\frac{1}{12}}\,\eta(\omega) = 1.
\]
